\chapter{Conclusion}\label{chap:conclusion}

We have introduced layout-agnostic and traversal-agnostic algorithm design approaches as practical solutions for enabling user-directed optimizations and performance portability of high-performance scientific applications, particularly those that are memory-bound. We have also demonstrated that these approaches simplify parallelization and distribution of applications on both CPU and GPU-based systems.

Noarr and Cellato serve as proof-of-concept implementations of abstractions supporting layout-agnostic and traversal-agnostic algorithm design in \cpp for general-purpose scientific and cellular automata applications, respectively. Both libraries are developed as header-only libraries, making them easy to integrate into existing applications. They are open-source and available on GitHub.

The two libraries have been evaluated on a variety of use cases for their respective domains, demonstrating they do not introduce significant overhead compared to hand-optimized implementations, while providing significant flexibility in terms of changing data layouts and traversal orders without modifying the core algorithm logic. Furthermore, Cellato provides domain-specific optimizations for cellular automata applications.

We have also explored the potential of using reasoning LLMs to automatically generate optimized code. Our findings show that LLMs benefit from user-directed guidance (tutoring) for specific optimization goals or design choices in cases where the algorithm is too complex for the LLM to optimize on its own. On the other hand, our findings also suggest that abstractions promoting some level of layout-agnostic and traversal-agnostic algorithm design in \cpp do not significantly benefit LLMs --- in many cases, hindering their ability to generate optimized code.

Open questions remain regarding sparse data structures in layout-agnostic and traversal-agnostic algorithm design, as well as other specific domains such as graph algorithms. It should be also noted that future \cpp features, such as reflection, may enable even more powerful and flexible abstractions for layout-agnostic and traversal-agnostic algorithm design. In the area of LLMs, further research may explore whether our findings generalize to newer, agentic LLMs, which can interact with tools and thus receive feedback on their code generation, enabling them to iteratively improve the generated code.

% Overall, our work demonstrates feasibility and benefits of layout-agnostic and traversal-agnostic algorithm design in high-performance scientific applications
The findings of this thesis should encourage further adoption of layout-agnostic and traversal-agnostic algorithm design approaches in scientific applications, as well as further research into their potential benefits and limitations. We hope that our work will contribute to the development of more flexible and portable high-performance applications in even more domains, and to the advancement of LLM usage for code generation in scientific computing.
